%% kapa: exam=1 round=23 year=2012 period=1 form=A subject=civil_law subject_ko=민법 seq=10 number=10
%% source: Q-Net 감정평가사 기출문제 (cst00309) artlSeq=5052578 file=raw/1008/5052578/2012년 제23회 감정평가사 제1차시험 기출문제(1교시 A형).pdf
%% conversion: kordoc 4.24.51 -> tools/gen1.py (auto); review=proofread-sonnet; flags=none
\begin{question}{10}
혼동에 관한 설명으로 옳지 않은 것은? (다툼이 있으면 판례에 의함)
\begin{choices}
  \item 乙 소유의 토지에 지상권을 취득한 甲이 그 지상권을 목적으로 하는 저당권을 丙에게 설정한 후 甲이 토지의 소유권을 취득하였다면, 甲의 지상권은 소멸하지 않는다.
  \item 乙 소유의 토지에 지상권을 취득한 甲이 그 지상권을 목적으로 하는 저당권을 丙에게 설정한 후에 甲이 그 저당권을 취득한다면, 그 저당권은 소멸한다.
  \item 甲 소유의 건물에 乙의 대항력 있는 임차권이 성립한 후 丙이 그 건물에 저당권을 취득하였다면, 후일 乙이 甲으로부터 그 건물을 매수하여도 임차권은 소멸하지 않는다.
  \item 丙 소유의 토지에 甲의 1순위저당권과 乙의 2순위저당권이 설정되어 있을 때, 甲이 그 토지를 매수하여 소유권을 취득하였고 甲의 피담보채권이 존속할 경우에는 乙의 저당권이 1순위가 된다.
  \item 근저당권자 甲이 근저당목적물인 건물을 매수한 후 그 소유권취득이 무효로 밝혀지면, 혼동으로 소멸하였던 근저당권은 부활한다.
\end{choices}
\end{question}
\answer{4}
